\section{The limiting electrical energy}\label{sec:energy}

We now identify the unnormalized energies of the primal and dual graphs.
The cutoffs of Proposition~\ref{nd:cutoffs} first give compactness for
harmonic functions and coercivity for arbitrary continuous limits.  The
locality of the extracted arrays then makes the limiting energy a
deterministic scalar multiple of the conformal Dirichlet integral.
An annular current--period argument will show that this scalar is exactly
one.  Thus the normalization is obtained from the network itself.

Throughout this section we work in an ordinary conformal chart, in a
coupled extraction with the conclusions of
Proposition~\ref{inp:array-interface} and Proposition~\ref{nd:cutoffs}.
Write $G_m,G_m^*$ for the primal and dual graphs and $\pi_m$ for their
surface projections.  All sets used below have fixed positive clearance
inside the chart before $m$ tends to infinity.  The energy
$\en_m(v;U)$ is the sum of $|v(y)-v(x)|^2$ over edges with both projected
endpoints in $U$; we use the same convention for compact sets.  Replacing
this convention by edge-location membership does not change an inner
limit, because the projected mesh tends to zero.  Uniform convergence
on a compact set always means convergence at every vertex over that set,
including every lift of a point.

\subsection{Circuits and harmonic compactness}

A cutoff on one graph gives a short level circuit on the other, where
``short'' is measured by the increments of the function under study.
This elementary coarea observation is the bridge between the cutoffs
and the analysis of functions.

\begin{lemma}[A separating circuit]\label{en:circuit}
Fix concentric disks with radii $0<r<R$ and with a neighborhood of the
closed larger disk inside an ordinary chart.  Given a primal function
$v_m$, there is a simple primal circuit $C_m$ in a slightly buffered
annulus which surrounds all lifts of the smaller disk and whose inside
contains no lift outside a slightly enlarged larger disk, such that
\begin{equation}\label{en:circuit-bound}
 \left(\sum_{e\in C_m}|\dd v_m(e)|\right)^2
 \le C\,\en_m(v_m;A^+).
\end{equation}
Here $A^+$ is a fixed buffered annulus and $C$ depends only on its aspect
and the chart constants.  The conclusion holds eventually for every
fixed annulus.  It also holds with the primal and dual graphs exchanged.
\end{lemma}

\begin{proof}
Take a dual cutoff $b_m$ which is one on the inner well and zero on the
outer well, with energy at most $C$.  These wells include all relevant
lifts, with fixed collars.  For $t\in(0,1)$ away from the finitely many
vertex values, color the primal faces according to whether their dual
vertices satisfy $b_m>t$.  The oriented boundary of the colored faces
splits into simple circuits.  A face appearing on both sides of an edge
contributes no boundary occurrence; other edge occurrences have their
usual multiplicities.

The boundary chain has winding difference one between an inner dual
vertex and an outer dual vertex, because their face colors are one and
zero.  After splitting the chain into simple circuits, at least one
circuit separates these two vertices.  The inner dual well is connected
through dual edges on which the cutoff is one; such a path cannot cross
any boundary edge, including at a primal articulation vertex.  Hence
every vertex of this inner well is on the same side of the chosen
circuit.  The connected outer well has the analogous property.  The
collars and vanishing overlay mesh put all primal inner lifts on the
inner side as well, and all lifts beyond the larger well on the outer
side.  This uses the all-lift joining statement of
Proposition~\ref{inp:topology} and remains valid when the full boundary
chain touches itself.  The selected circuit lies in the change annulus
and has variation at most that of the entire boundary.
The coarea identity and Cauchy--Schwarz give
\[
 \int_0^1\sum_{e\in\partial\{b_m>t\}}|\dd v_m(e)|\,\dd t
 =\sum_e |\dd v_m(e)|\,|\dd b_m(e^*)|
 \le \en_m(v_m;A^+)^{1/2}\en_m^*(b_m)^{1/2}.
\]
Choose a level at most this average.  This proves
\eqref{en:circuit-bound}, with an arbitrarily small slack if necessary.
The same argument applies after exchanging the graphs.
\end{proof}

\begin{lemma}[Harmonic compactness]\label{en:harmonic-compactness}
Let $U$ be a connected open subset of an ordinary chart.  Suppose that
$v_m$ is graph-harmonic at every vertex whose star projects into $U$,
and that its energies are bounded on compactly contained neighborhoods.
Then $v_m$ is macroscopically equicontinuous on every compact subset of
$U$: for $K\Subset U$,
\[
 \lim_{\delta\downarrow0}\limsup_{m\to\infty}
 \sup_{\substack{\pi_m x,\pi_m y\in K\\
                  |\pi_m x-\pi_m y|<\delta}}
                  |v_m(x)-v_m(y)|=0.
\]
If $x_m$ is a vertex sequence with $\pi_mx_m\to z_0\in U$ and
$v_m(x_m)$ is bounded, the functions have locally uniformly convergent
subsequences.  The assertions hold for both graphs and for local
branches of potentials, with a normalization on each initial branch.
\end{lemma}

\begin{proof}
Fix a disk with clearance inside $U$ and an energy bound $E$ on a larger
disk.  Around its center place $J$ disjoint annuli of one fixed aspect,
all inside a prescribed small neighborhood.  Choose all radii before
taking $m$ large.  Their energy sum is at most $E$, so
Lemma~\ref{en:circuit} supplies a surrounding circuit on which the
oscillation is at most $(CE/J)^{1/2}$.  The discrete maximum principle
bounds the oscillation everywhere inside by the same quantity.  Taking
$J$ large, and then using finitely many such neighborhoods, proves
equicontinuity.  No assertion is needed at a radius depending on $m$.

The same construction at a fixed, nonvanishing scale bounds oscillation
on compact balls.  Overlapping chains of balls propagate an additive
normalization throughout a connected compact subset.  Vertex density
and the small-image joining property identify the limits on overlaps.
The usual finite-net proof of Arzel\`a--Ascoli now gives local uniform
compactness.  A local potential branch is an ordinary harmonic function
on each of its branch disks, so the proof also applies to branches.
\end{proof}

\subsection{The uniform relaxation and its energy measures}

The extracted variational tests let us define a single limiting energy
for every continuous target.  This definition does not assume a
continuum Dirichlet form.

For $u\in C(U)$, $K\Subset U$ compact, and $\delta>0$, let
\[
 M_m(u,K,\delta)=
 \inf\left\{\en_m(v;K):
  |v(x)-u(\pi_m x)|\le\delta
  \text{ for every vertex }x\text{ over }K\right\}.
\]
Define the \emph{inner relaxed energy} by
\begin{equation}\label{en:relaxation}
 F(u,U)=\sup_{K\Subset U}\sup_{\delta>0}
                  \liminf_{m\to\infty}M_m(u,K,\delta).
\end{equation}
One may use a countable exhaustion of compact finite unions of boxes
and let $\delta$ decrease through positive rationals.  Define $F^*$ by
the same formula on the dual graph.

\begin{lemma}[Uniform variational limits]\label{en:envelope}
The value in \eqref{en:relaxation} is unchanged if $\liminf$ is replaced
by $\limsup$.  It is lower semicontinuous for locally uniform convergence
and inner regular in $U$.  Moreover:
\begin{enumerate}[label=\textup{(\roman*)}]
\item Every locally uniformly convergent sequence $v_m\to u$ has local
energy lower bounds given by $F(u,\cdot)$, also along further
subsequences.
\item If $F(u,U)<\infty$, there are compact sets $K_m$ exhausting $U$
and functions $v_m$ on their graph neighborhoods such that $v_m\to u$
locally uniformly and
$\en_m(v_m;K_m)\le F(u,U)+o(1)$.
\end{enumerate}
These statements hold simultaneously for all continuous targets.  The
values $F(u,U)$ are locally readable from the arrays with arbitrarily
small open buffers.
\end{lemma}

\begin{proof}
Given $K\Subset K'\Subset U$ and $0<\delta'<\delta$, insert between the
two tests a finite selector test from Lemma~\ref{inp:sandwich}.  Cover
the relevant visit labels by finitely many selectors on which the target
varies by less than a fixed fraction of $\delta-\delta'$, and use
intermediate rational value constraints.  Completeness of the selected
graph preserves every edge incidence.  Restriction of an admissible
function for the larger test is admissible for the intermediate one;
restriction from the intermediate test is admissible for the smaller
one.  The intermediate infimum has a retained limit.  Consequently,
\[
 \limsup_m M_m(u,K,\delta)
 \le \liminf_m M_m(u,K',\delta').
\]
Taking the two inner suprema proves the assertion about $\liminf$ and
$\limsup$, including infinite values.

The same comparison proves the lower bound in (i).  For (ii), choose
almost minimizers on increasing compact sets with decreasing
tolerances, and take a diagonal sequence using the $\limsup$ formula.
Tolerance slack also proves lower semicontinuity, directly from
\eqref{en:relaxation}.  Exhausting $U$ proves inner regularity.  The
selector library is countable; approximation of a continuous target
uses its modulus of continuity on each compact, rather than a new
exceptional event.  This gives the simultaneous assertion.  The
locality and measurability claims follow from the same selector
sandwich, with its spatial buffers.
\end{proof}

We call the sequences in Lemma~\ref{en:envelope}(ii) recovery sequences.
The next estimate explains why such sequences can be joined on overlaps.
For an oriented edge $xy$, put
$\overline v(xy)=(v(x)+v(y))/2$.

\begin{lemma}[Gluing estimate]\label{en:fundamental}
Let $v_m,w_m$ converge uniformly to the same continuous function on a
fixed transition region.  Suppose $0\le b_m\le1$ has bounded energy,
is supported where $v_m$ is defined, and equals one wherever $w_m$ is
not used.  Then $z_m=b_mv_m+(1-b_m)w_m$ has, on the gluing region,
\begin{align}\label{en:gluing-estimate}
 \en_m(z_m)
 &\le(1+\varepsilon)\sum_e
 \left(\overline b_m(e)|\dd v_m(e)|^2
 +(1-\overline b_m(e))|\dd w_m(e)|^2\right)\notag\\
 &\quad +(1+\varepsilon^{-1})
       \|v_m-w_m\|_{\infty,\mathrm{tr}}^2\en_m(b_m),
 \qquad \varepsilon>0.
\end{align}
Finite partitions of unity have the corresponding estimate, with an
error tending to zero when the local approximations agree uniformly on
overlaps.  Bounded-energy gluing weights exist for every fixed finite
open cover of a compact set.
\end{lemma}

\begin{proof}
The exact difference identity is
\[
 \dd z_m=\overline b_m\,\dd v_m
        +(1-\overline b_m)\,\dd w_m
        +\overline{v_m-w_m}\,\dd b_m.
\]
Apply convexity of the square to the first two terms, followed by
$|a+c|^2\le(1+\varepsilon)|a|^2+
(1+\varepsilon^{-1})|c|^2$.  Only the transition region contributes to
the last term.  For finitely many weights the same identity follows by
summing and using $\sum_i\dd b_{i,m}=0$.

To construct weights, choose finitely many disks whose smaller disks
cover the compact set and whose larger disks stay in the prescribed
members of the cover.  Use Proposition~\ref{nd:cutoffs}, sum the
resulting cutoffs, and normalize where their sum is at least one.
Clipping and division on this fixed range preserve bounded energy.
Extra cutoffs with collars extend the construction to a neighborhood
of the compact set.  All these choices precede the limit in $m$.
\end{proof}

\begin{proposition}[Quadratic energy measures]\label{en:measure}
For each open $U$, the functional $F(\cdot,U)$ is quadratic on its
finite-energy space and vanishes on constants.  If $u$ has locally
finite relaxed energy, then
\[
                    U\longmapsto F(u,U)
\]
is the restriction to open sets of a Radon measure.  For smooth $u$,
on a smaller ordinary chart,
\begin{equation}\label{en:smooth-upper}
                  F(u,U)\le C\int_U|\nabla u|^2\,\dd z.
\end{equation}
The same statements hold for $F^*$.
\end{proposition}

\begin{proof}
Recovery sequences for $u$ and $v$, the discrete parallelogram identity,
and the energy lower bound give
\[
 F(u+v,U)+F(u-v,U)\le2F(u,U)+2F(v,U).
\]
Applying this inequality to $u+v$ and $u-v$ gives the reverse inequality.
Homogeneity follows directly from the tolerance definition.  Constant
addition leaves every edge increment unchanged.

For fixed $u$, restriction gives monotonicity and superadditivity on
disjoint open sets.  To prove subadditivity on $U\cup V$, first fix a
compact subset of this union.  Recover $u$ in $U$ and $V$, and join the
recoveries with Lemma~\ref{en:fundamental}; the transition is compactly
contained in $U\cap V$.  Let $m\to\infty$ and then
$\varepsilon\downarrow0$ in \eqref{en:gluing-estimate}.  Exhausting the
compact gives
\[
                   F(u,U\cup V)\le F(u,U)+F(u,V).
\]
Finite subcovers and inner regularity imply countable subadditivity.
For completeness, these properties give the required measure criterion
as follows.  Define
$\nu(A)=\inf\{F(u,O):A\subset O,\ O\text{ open}\}$.
This is an outer measure by countable subadditivity.  If two sets have
positive distance, they have disjoint open neighborhoods; disjoint-open
superadditivity shows that $\nu$ is additive on such sets.  Thus it is a
metric outer measure and all Borel sets are measurable.  On an open set
$O$, monotonicity gives $\nu(O)=F(u,O)$.  Local finiteness and inner
regularity make this Borel measure Radon.

It remains to prove \eqref{en:smooth-upper}; this will also establish
local finiteness for smooth targets.  On a compact set inside $U$, use a
square grid of spacing $r$ and a bounded-overlap family of disk cutoffs
with radii comparable to $r$.  Normalize these cutoffs to weights
$b_{z,m}$ which sum to one near the compact, and set
\[
                         v_m=\sum_z b_{z,m}u(z).
\]
For fixed $r$ this approximates $u$ within its modulus of continuity at
scale $Cr$.  On each edge the weight increments sum to zero.  Subtracting
one nearby node value and using bounded overlap bounds the energy by
\[
 C\sum_z r^2\sup_{B(z,C'r)}|\nabla u|^2.
\]
Indeed each unnormalized cutoff has bounded energy, and normalization
has the same bound because the denominator is bounded below and only
boundedly many supports meet an edge.  First take $m$ large for this
finite collection, and then let $r\downarrow0$.  The displayed sum
converges to at most a constant times the gradient integral over the
chosen region.  Exhaustion inside $U$ proves
\eqref{en:smooth-upper}.  The arguments are unchanged on the dual.
\end{proof}

\subsection{Integral representation and coercivity}

The energy measures of the two coordinate functions determine the
whole functional.  Coercivity is a separate step: it follows from
intersecting level circuits and does not assume geometric regularity of
the mesh.

\begin{proposition}[Integral representation]\label{en:representation}
In every ordinary chart there is a measurable symmetric matrix field
$A$, locally bounded and uniformly positive definite on smaller charts,
such that, for every open $U$ and every $u\in C(U)$,
\begin{equation}\label{en:matrix-form}
 F(u,U)=
 \begin{cases}
 \displaystyle\int_U\nabla u^{\mathsf T}A\nabla u\,\dd z,
                &u\in W^{1,2}_{\mathrm{loc}}(U),\\[4pt]
 +\infty,       &u\notin W^{1,2}_{\mathrm{loc}}(U).
 \end{cases}
\end{equation}
The integral is allowed to be infinite.  There is an analogous field
$A^*$ for the dual functional.
\end{proposition}

\begin{proof}
We first identify the energy of smooth functions.  Polarize the Radon
measures from Proposition~\ref{en:measure}.  The measures of the linear
functions $\ell_1(z)=\Re z$, $\ell_2(z)=\Im z$ and their sums are
absolutely continuous by \eqref{en:smooth-upper}.  Define the entries of
$A$ by their densities:
\[
 A_{ii}\,\dd z=F(\ell_i,\dd z),\qquad
 2A_{12}\,\dd z=
 F(\ell_1+\ell_2,\dd z)-F(\ell_1,\dd z)-F(\ell_2,\dd z).
\]
Quadraticity and the upper bound imply that $A$ is symmetric,
nonnegative, and locally bounded.  At a Lebesgue point $z$ of these
densities, compare a smooth $u$ with its affine first-order Taylor
polynomial $p_z$.  The seminorm inequality and
\eqref{en:smooth-upper} give
\[
 \left|F(u,B(z,r))^{1/2}-F(p_z,B(z,r))^{1/2}\right|
 \le C^{1/2}|B(z,r)|^{1/2}
          \sup_{B(z,r)}|\nabla u-\nabla u(z)|.
\]
Divide by $|B(z,r)|^{1/2}$ and let $r\downarrow0$.  Differentiation of
Radon measures yields \eqref{en:matrix-form} for smooth $u$.

We next prove the converse bound, including weak differentiability:
\begin{equation}\label{en:coercivity}
                 \int_U|\nabla u|^2\,\dd z\le C F(u,U).
\end{equation}
Suppose the right side is finite.  Put a square grid of fixed small
spacing $r$ in a compactly contained region of $U$.  At every grid node
$x$, apply Lemma~\ref{en:circuit} to a recovery sequence, using an
annulus with inner and outer radii $0.7r$ and $0.8r$, and choose a value
$a_{x,m}$ at a vertex on the resulting circuit.  Buffers are small fixed
fractions of $r$.

Circuits at neighboring grid nodes share a vertex.  Here is the
necessary topological verification.  A point near the midpoint of the
two centers is inside both inner disks.  Points on the two opposite
sides belong to exactly one inner disk and lie beyond the other outer
disk; a distant point lies outside both.  The all-lift separation rule
therefore gives a common interior lift, a common exterior lift, and
lifts with each of the two opposite memberships.  Two disjoint simple
circuits on the sphere with these memberships could be neither nested
nor disjoint in their interiors, a contradiction.  Since both circuits
belong to the primal graph, their intersection contains a vertex.
Consequently,
\[
 |a_{x,m}-a_{y,m}|^2
 \le 2\left(\sum_{e\in C_{x,m}}|\dd v_m(e)|\right)^2
      +2\left(\sum_{e\in C_{y,m}}|\dd v_m(e)|\right)^2
\]
for neighbors $x,y$.  Write $B_r$ for the union of the buffered grid
annuli.  Bounded overlap and \eqref{en:circuit-bound} give
\[
 \sum_{x\sim y}|a_{x,m}-a_{y,m}|^2
 \le C\en_m(v_m;B_r)
 \le C\bigl(F(u,U)+o(1)\bigr).
\]
For the last inequality, this fixed finite union lies in the recovery
set $K_m$ for all sufficiently large $m$.
The piecewise affine interpolation of the $a_{x,m}$ on grid triangles
has Dirichlet energy bounded by this sum times an absolute constant.
Take $m$ large after fixing the grid, and then let $r\downarrow0$.
Continuity of $u$ and uniform recovery imply that these interpolations
converge locally uniformly to $u$.  Weak Sobolev compactness proves
\eqref{en:coercivity} on the compact region; exhaustion proves it on
$U$.  Applying it to smooth affine functions on arbitrary disks gives
$A\ge cI$ locally.

Finally, let $u\in C(U)\cap W^{1,2}_{\mathrm{loc}}(U)$.  On smaller
regions, smoothing approximates $u$ both uniformly and in $W^{1,2}$.
Lower semicontinuity and the smooth upper bound first extend that upper
bound to such continuous Sobolev functions.  Apply the extended bound
to the difference between $u$ and its smooth approximations.  Their
$F$-seminorm difference tends to zero, so the smooth formula passes to
$u$.  Exhaustion gives \eqref{en:matrix-form}, including infinite
integrals.  If local Sobolev regularity fails, \eqref{en:coercivity}
forces infinite energy.  All steps apply equally to $F^*$.
\end{proof}

\begin{lemma}[The deterministic scalar]\label{en:scalar}
In each coupled extraction there is a deterministic
$\lambda\in(0,\infty)$ such that
\begin{equation}\label{en:scalar-form}
                    A=A^*=\lambda I
\end{equation}
almost everywhere in all ordinary charts.  At this stage $\lambda$
may depend on the extracted rate sequence.
\end{lemma}

\begin{proof}
The derivatives of the energy measures of linear functions are point
germs of the arrays: use the tolerance definition on successively
smaller disks, with arbitrarily small open buffers.  Choose measurable
$\limsup$ versions and a fixed default on the null set where
Lebesgue differentiation fails.  Lemma~\ref{nd:germs} applies to these
matrix-valued readouts (or to a bounded injective encoding of their
entries), so $A$ and $A^*$ are deterministic almost everywhere.
Translation covariance makes each field a constant matrix on overlapping
ordinary charts.  Under a rotation, the matrix transforms by
$A\mapsto RAR^{\mathsf T}$.  Rotation covariance, in particular a
quarter-turn, therefore makes each symmetric matrix scalar.  These
uses of covariance are precisely those in
Remark~\ref{nd:germ-symmetries}, and compare local laws with their actual
field transformations.  Primal--dual reflection
(Lemma~\ref{inp:reflection}) equates the two scalar values.  The upper
bound and coercivity give $0<\lambda<\infty$.
\end{proof}

\subsection{Annular reciprocity fixes the scalar}

We have proved existence and isotropy of the limiting energy.  To
identify its normalization, compare a condenser potential with a
potential whose value increases by one on going around an annulus.
The first comparison rotates a discrete current to a dual potential;
the reverse comparison recovers an angular potential with its period
\emph{exactly} preserved.  Both comparisons use only interior graph
pieces with collars.  The current--potential correspondence is the
standard planar-network duality; see \cite[Section~9.2]{LP}.
Its use to determine a scalar effective conductivity has a classical
precedent in Dykhne's planar self-duality argument \cite{Dykhne}.
Here we give both inequalities, including the recovery needed on our
irregular graph pieces.

Let
\[
 A(r,R)=\{z:r<|z|<R\},\qquad
 C_0=\frac{2\pi}{\log(R/r)},
\]
where the closed annulus and a neighborhood of it lie inside an ordinary
chart.  Thus $C_0$ is the ordinary Dirichlet capacity between its two
boundary circles.  The radial minimizer is
$\log(|z|/r)/\log(R/r)$.  A locally Sobolev potential with additive
period $P$ around the annulus has energy at least
\begin{equation}\label{en:period-bound}
             \int_{A(r,R)}|\nabla v|^2\,\dd z
                    \ge\frac{P^2}{C_0}.
\end{equation}
Indeed, on almost every circle,
$\int_0^{2\pi}\partial_\theta v\,\dd\theta=P$; apply
Cauchy--Schwarz and integrate $\dd s/s$.  Equality holds for
$v=P\theta/(2\pi)$, interpreted in local branches.

\begin{lemma}[Condenser convergence]\label{en:condenser}
Prescribe $u_m=0$ on every primal vertex over $|z|\le r$ and $u_m=1$
on every vertex over $|z|\ge R$, using a protected cap outside the chart.
Let $C_m$ be the minimum primal energy.  Then
\begin{equation}\label{en:condenser-limit}
                             C_m\longrightarrow\lambda C_0.
\end{equation}
The minimizers converge uniformly, including at the two plate
boundaries, to the round radial condenser potential.
\end{lemma}

\begin{proof}
All nonzero increments of the minimizer project into any fixed collar
of the closed annulus eventually.  Indeed, both endpoints of an edge
strictly on one plate have its prescribed value, and the overlay mesh
vanishes.  Thus the cap contributes no uncontrolled energy term.

For the upper bound, approximate the radial potential by smooth
functions which are constant on slightly wider plate collars and whose
energies tend to $C_0$.  Recover them uniformly on a neighborhood of the
annulus.  Bounded-energy weights in the flat collars paste in the exact
constants at a vanishing energy cost, by
Lemma~\ref{en:fundamental}.  Lemma~\ref{en:scalar} gives
$\limsup C_m\le\lambda C_0$.

The minimizer lies in $[0,1]$ and is harmonic off the plates.  Its bounded
energy and Lemma~\ref{en:harmonic-compactness} give compactness in the
open annulus.  At a plate-boundary point, use the stacks of rings from
that lemma, small enough to avoid the other plate.  Every surrounding
circuit meets a path lying strictly on the plate side: choose a tube
there from the inner well to outside the ring and lift it by
Proposition~\ref{inp:topology}.  Thus the circuit contains a vertex
with the prescribed value.  A circuit of small variation is uniformly
close to that value.  Inside it all remaining vertices are harmonic or
have that same prescribed value, so the maximum principle proves
continuity uniformly up to the plate.  Fixed stacks followed by finite
covers give the assertion simultaneously along both boundaries.

Any uniform subsequential limit has the condenser boundary values.
The variational lower bound and Lemma~\ref{en:scalar} give energy at
least $\lambda C_0$.  This proves \eqref{en:condenser-limit}; equality
in the round Dirichlet principle identifies the limit and gives the
last assertion.
\end{proof}

\begin{figure}[tb]
\centering
\begin{tikzpicture}[scale=1.05,>=Stealth]
 \fill[gray!10,even odd rule] (0,0) circle (1.45) (0,0) circle (.55);
 \draw (0,0) circle (1.45);
 \draw (0,0) circle (.55);
 \node[font=\small] at (0,0) {$u_m=0$};
 \node[anchor=east] at (-1.52,0) {$u_m=1$};
 \foreach \a in {40,140,220,320}{
   \draw[blue!65!black,->,thick]
       (\a:.64) -- (\a:1.29);
 }
 \draw[red!65!black,->,thick] (18:1.04)
       arc[start angle=18,end angle=153,radius=1.04];
 \node[red!65!black,anchor=south] at (0,1.54)
       {dual period $C_m$};
 \node[blue!65!black,anchor=north] at (0,-1.6)
       {primal current $\dd u_m$};
 \draw[->,thick] (1.9,.6) -- (3.1,.6);
 \node[align=center] at (2.5,.05) {unit-period\\recovery};
 \begin{scope}[xshift=5cm]
  \fill[gray!10,even odd rule] (0,0) circle (1.45) (0,0) circle (.55);
  \draw (0,0) circle (1.45);
  \draw (0,0) circle (.55);
  \draw[red!65!black,->,thick] (18:1.04)
       arc[start angle=18,end angle=330,radius=1.04];
  \node at (0,0) {$\alpha_m$};
  \node[red!65!black,anchor=south] at (0,1.54)
       {exact period $1$};
  \node[anchor=north] at (0,-1.6)
       {$|\langle\dd u_m,\alpha_m\rangle|=1$};
 \end{scope}
\end{tikzpicture}
\caption{The two annular comparisons.  Rotating the condenser current
gives a dual potential with period $C_m$ and hence $\lambda\le1$.
An exact unit-period dual recovery pairs with the unit condenser jump
to give $\lambda\ge1$.  This is a schematic of the test annulus;
the projected discrete graph need not resemble a regular mesh.}
\label{en:fig-reciprocity}
\end{figure}

\begin{lemma}[Rotation of a current]\label{en:current-period}
Rotate the primal current $\dd u_m$ from
Lemma~\ref{en:condenser} to the dual edges, with one consistent crossing
orientation.  In the annulus this is a closed dual cochain with harmonic
local potential branches and period $\pm C_m$.  The branches have
subsequential locally uniform limits with period $\pm\lambda C_0$ and
with total dual relaxed energy at most $\lambda C_0$ over one turn.
\end{lemma}

\begin{proof}
Closedness around a dual face is the primal harmonic equation.
Harmonicity of a potential branch at a dual vertex is the telescoping
sum of $\dd u_m$ around the corresponding primal face.  These identities
count all incidences and are valid also for loops and repeated edges.
The squared dual increments sum to at most $C_m$ on every interior
annular piece.

We justify the existence of branches without assuming regular discrete
annular boundaries.  In a disk with clearance inside the annulus,
arbitrary dual lifts in a smaller disk are connected within the larger
one.  Loops used to integrate the cochain can be filled inside a further
buffer, by Proposition~\ref{inp:topology}; every primal vertex there is
harmonic.  Cellular closedness therefore makes the integral independent
of the chosen path.  On connected overlaps the resulting potentials
differ by constants.

A dual loop with projected winding one can be built along fixed tubes
and closed using small-image joining.  The integer-winding statement of
Proposition~\ref{inp:topology} makes its winding upstairs one, up to the
fixed orientation sign.  Its current integral is the total current
entering one plate, since all intervening primal vertices are harmonic.
Discrete summation by parts gives total current $C_m$: the potential
difference is one and its energy is $C_m$.  Hence the period is
$\pm C_m$.

Normalize branches on a finite overlapping disk cover of any compact
subannulus, propagating constants through chosen overlaps.
In particular, fix a cyclic chain of disks around one positively
oriented circle.  Continue a normalized branch along this chain,
matching on each connected overlap; on returning to the first disk,
the continued branch differs from the original by $\pm C_m$.
Lemma~\ref{en:harmonic-compactness} gives local uniform subsequential
limits.  Their transition constants converge, so their period is
$\pm\lambda C_0$.  Exhausting the annulus yields branches everywhere.
For the energy bound, use disjoint boxes lying in branch disks.  Their
discrete edge energies do not overlap, and their sum is at most $C_m$.
Apply the local lower bound and then exhaust the annulus up to sets of
Lebesgue measure zero.  This gives the asserted bound over one turn.
\end{proof}

Combining Lemmas~\ref{en:scalar} and \ref{en:current-period} with
\eqref{en:period-bound} gives
\begin{equation}\label{en:lambda-upper}
        \lambda C_0\ge
        \lambda\frac{(\lambda C_0)^2}{C_0},
        \qquad\text{hence }\lambda\le1.
\end{equation}
For the reverse inequality, angular recovery must retain its period
before taking any limit.  We record first the measure observation that
makes the gluing sharp.

\begin{lemma}[Energy measures of recovery sequences]\label{en:recovery-measures}
Let $V$ be a bounded patch and let $p$ be smooth on a neighborhood of its
closure.  Choose a recovery sequence on compact sets exhausting $V$
whose total energy tends to $F^*(p,V)$.  Count only edges with both
endpoints in these recovery sets, and put their squared increments
at any of the corresponding projected edge locations.  These measures
converge vaguely in $V$ to $\lambda|\nabla p|^2\,\dd z$.  In particular,
energies weighted by a compactly supported continuous function converge
to the corresponding weighted integral.
\end{lemma}

\begin{proof}
The total masses are bounded.  Let $\nu$ be a subsequential vague limit.
To check its local lower bound, take open sets $W\Subset O\Subset V$
and a compact set $K\subset O$ containing $W$ with a collar.  The energy
tests inside $W$ are eventually counted by the edge measure $\nu_m(K)$.
Lemma~\ref{en:envelope} and compact-set Portmanteau give
\[
 F^*(p,W)\le\liminf_m\nu_m(K)
          \le\limsup_m\nu_m(K)\le\nu(K)\le\nu(O).
\]
Exhaust $O$ by such $W$.  Thus $\nu$ dominates
$\lambda|\nabla p|^2\,\dd z$.  Its total mass is at most the limiting
total recovery energy, which is the total mass of that same measure.
Equality follows.  Every subsequential vague limit is consequently the
stated measure, proving the claim.
\end{proof}

\begin{lemma}[Recovery with exact unit period]\label{en:angular-recovery}
For every sufficiently small $\eta>0$, let
$A_\eta=A(r-\eta,R+\eta)$.  On the dual edges in a neighborhood of the
closed annulus there are closed cochains $\alpha_m$ with exact period
one and
\begin{equation}\label{en:angular-energy}
 \limsup_{m\to\infty}\sum_{e^*\subset A_\eta}|\alpha_m(e^*)|^2
 \le\lambda\left(\frac1{C_0}+o_\eta(1)\right),
             \qquad o_\eta(1)\longrightarrow0.
\end{equation}
The cochains are closed at every cell whose star lies in this
neighborhood; the region can be chosen to contain every edge where
$\dd u_m\ne0$ for all large $m$.
\end{lemma}

\begin{proof}
Take finitely many disks covering a slightly larger closed annulus,
each avoiding the origin, with angle branches
$p_i=\theta_i/(2\pi)$.  Recover $p_i$ uniformly on larger branch disks.
Lemma~\ref{en:recovery-measures} applies to these recoveries, denoted
$v_{i,m}$.  Take a smooth partition of unity $b_i$ with supports compactly
inside the branch disks.  Recover its weights, clip them to be
nonnegative, cut them off in fixed support collars, and normalize their
sum.  This gives $b_{i,m}\to b_i$ uniformly, bounded energies, supports
inside the branch disks, and $\sum_i b_{i,m}=1$ near the annulus.

There is a useful canonical cochain even though its own energy need
not be bounded.  On an oriented dual edge, let
\[
 a_m(e^*)=\frac1{2\pi}\int_{\pi_m(e^*)}\dd\arg z.
\]
Here the integral means the endpoint increment of a continuous argument
lift along the projected edge.  Every sufficiently small edge image
lies in a branch disk, so this is well defined.  Vanishing face mesh
makes $a_m$ closed around every cell under consideration, and its
period around a loop is that loop's projected integer winding.

On the support of $b_{i,m}$ the error
$v_{i,m}-p_i\circ\pi_m$ is a single-valued function tending uniformly
to zero.  Define the globally single-valued error and cochain by
\begin{equation}\label{en:angular-error}
 q_m=\sum_i b_{i,m}(v_{i,m}-p_i\circ\pi_m),
                 \qquad \alpha_m=a_m+\dd q_m.
\end{equation}
The support collars make this definition independent of how terms are
extended by zero.  Since a gradient has zero periods, $\alpha_m$ is
closed and has period exactly one, in the projected orientation.
This conclusion is an identity at each finite resolution.

To bound its energy, apply the discrete product identity on each edge.
On a branch disk, $\dd(p_i\circ\pi_m)=a_m$.  Since the weights sum to
one, the terms containing $a_m$ cancel to give
\[
 \alpha_m(e^*)=
 \sum_i\overline b_{i,m}(e^*)\,\dd v_{i,m}(e^*)
 +\sum_i
 \overline{v_{i,m}-p_i\circ\pi_m}(e^*)\,\dd b_{i,m}(e^*).
\]
The second sum tends to zero in edge $\ell^2$, by uniform convergence
of the errors and bounded energy of the finitely many weights.
Convexity bounds the squared norm of the first sum by the weighted sum
of local recovery energies.  Lemma~\ref{en:recovery-measures} and
$b_{i,m}\to b_i$ turn its limit into
\[
 \lambda\sum_i\int_{A_\eta} b_i|\nabla p_i|^2\,\dd z
 =\lambda\int_{A_\eta}\left|\nabla\frac{\theta}{2\pi}\right|^2\,\dd z
 =\frac{\lambda}{2\pi}\log\frac{R+\eta}{r-\eta}.
\]
One may compute first on a slightly larger closed annulus inside the
cover; its boundary circles have zero limiting energy measure.
The edge averages $\overline b_{i,m}$ converge uniformly to $b_i$ at
the chosen edge locations, by vanishing mesh.  This justifies both
the weighted passage to the limit and restriction to the annulus.
This proves \eqref{en:angular-energy}.  Finally, $u_m$ is constant on
each plate, so its nonzero increments lie in $A_\eta$ eventually by
vanishing edge mesh.  Choosing additional fixed collars gives the
assertion about stars.
\end{proof}

\begin{proposition}[Exact conductance normalization]\label{en:reciprocity}
The scalar in Lemma~\ref{en:scalar} equals one.
\end{proposition}

\begin{proof}
Only $\lambda\ge1$ remains after \eqref{en:lambda-upper}.
Pair $\dd u_m$ with the cochain $\alpha_m$ of
Lemma~\ref{en:angular-recovery}, matching each primal edge with its
crossing dual edge and using the consistent crossing sign.  Then
\begin{equation}\label{en:pairing}
                         \left|\sum_e\dd u_m(e)\alpha_m(e^*)\right|=1.
\end{equation}
To see this at the discrete level, replace $u_m$ by the indicator of the
outer side of an intermediate circle.  Their difference is supported
strictly inside the region of closedness of $\alpha_m$.  Cellular
summation by parts therefore leaves the pairing unchanged.  The
indicator's gradient pairs with $\alpha_m$ along the oriented dual
boundary chain separating the two wells.  This chain has total winding
one up to sign by Proposition~\ref{inp:topology}; the cochain has exact
unit period by \eqref{en:angular-error}.  This proves
\eqref{en:pairing}.  The argument uses the boundary chain with its
multiplicities and does not require a simple discrete plate boundary.

Cauchy--Schwarz, \eqref{en:condenser-limit}, and
\eqref{en:angular-energy} now give
\[
 1\le C_m\sum_{e^*\subset A_\eta}|\alpha_m(e^*)|^2
 \quad\Longrightarrow\quad
 1\le\lambda C_0\,\lambda\left(C_0^{-1}+o_\eta(1)\right).
\]
Let $\eta\downarrow0$.  Thus $1\le\lambda^2$, and
\eqref{en:lambda-upper} proves $\lambda=1$.
\end{proof}

\begin{theorem}[Uniform conformal energy limit]\label{en:variational}
In every coupled extraction and every ordinary chart, both unnormalized
network energies have the same inner relaxed limit:
\begin{equation}\label{en:final-energy}
                   F(u,U)=F^*(u,U)=\int_U|\nabla u|^2\,\dd z
\end{equation}
for every $u\in C(U)\cap W^{1,2}_{\mathrm{loc}}(U)$, with infinite value
when the integral diverges or local Sobolev regularity fails.  The local
uniform lower bounds and recovery conclusions of
Lemma~\ref{en:envelope}, the gluing estimate, and harmonic compactness
hold simultaneously.  In particular, the limiting energy has no
subsequence-dependent multiplicative constant.
\end{theorem}

\begin{proof}
Combine Proposition~\ref{en:representation}, Lemma~\ref{en:scalar}, and
Proposition~\ref{en:reciprocity}.  The preceding construction used the
same countable buffered tests throughout, so its simultaneous and
locality assertions persist.  All arguments apply in every fixed
quantum-unit offset of the extraction.  Their compatibility is the
transport statement of Proposition~\ref{inp:array-interface}.
\end{proof}
